\begin{pastexam}[label={px:gw3-104-ds-03}]{關務三等資訊處理 104 年 第三題}
（申論；20 分）

佇列（Queue）結構的插入（Insert）和刪除（Delete）演算法如下：

{\small
\noindent const int N=10; int Rear=0, Front=0;\\[2pt]
\begin{minipage}[t]{0.47\linewidth}
\begin{tabbing}
\hspace*{1.2em}\=\hspace*{1.5em}\=\hspace*{1.5em}\=\kill
void Insert(char item, char Queue[])\\
\{ \> if (Rear==N-1)\\
\>\> cout$<<$``Queue Is Full'';\\
\> else\\
\> \{ \> Rear=Rear+1;\\
\>\> Queue[Rear]=item;\\
\> \}\\
\}
\end{tabbing}
\end{minipage}\hfill
\begin{minipage}[t]{0.47\linewidth}
\begin{tabbing}
\hspace*{1.2em}\=\hspace*{1.5em}\=\hspace*{1.5em}\=\kill
void Delete(char item, char Queue[])\\
\{ \> if (Front==Rear)\\
\>\> cout$<<$``Queue Is Empty'';\\
\> else\\
\> \{ \> Front =Front + 1;\\
\>\> item = Queue[Front];\\
\> \}\\
\}
\end{tabbing}
\end{minipage}
}
\begin{itemize}[leftmargin=2.2em, nosep]
\item[㈠] 請問上述演算法的佇列結構，會有什麼問題存在？（5 分）
\item[㈡] 可用什麼資料結構解決？（5 分）
\item[㈢] 承上之資料結構，請寫出插入（Insert）和刪除（Delete）演算法。（10 分）
\end{itemize}
\end{pastexam}

\begin{pastsoln}
答案：㈠ 線性佇列的空間不會回收：Rear 到了 N-1 就說滿了，即使前面已經刪掉很多、空出位置（假溢位）。 ㈡ 環狀佇列
（circular queue）。 ㈢ 程式在下面。

\textbf{㈠ 問題。} 這是線性佇列：Front 指向第一個元素的前一格，Rear 指向最後一個元素，兩者都只會往後移。
\begin{itemize}
\item \textbf{假溢位}：Insert 只看 Rear 是否到了 N-1。 例如連續插入 $9$ 個（Rear $= 9$）再刪除 $9$ 個（Front $= 9$），
佇列是空的，Insert 卻仍然說「Queue Is Full」。 刪除後空出的 Queue[0..Front] 永遠用不到。
\item \textbf{位置 0 永遠不用}：Rear 先加一再存，第一個元素放在 Queue[1]，Queue[0] 浪費掉，最多只能同時存 N-1 $= 9$ 個。
\item 要避免假溢位，只能在 Rear 到底時把所有元素往前搬，每次 $O(n)$，很慢。
\item （另外，Delete 的 \texttt{item} 是傳值參數，取到的值傳不回呼叫端，應改成傳參考 \texttt{char \&item} 或用回傳值；
下面的寫法一併修正。）
\end{itemize}

\textbf{㈡} 用\textbf{環狀佇列}：把陣列想成頭尾相接的環，索引用 $\bmod N$ 前進，Rear 到了 N-1 之後繞回 0，前面空出的位置
可以重複使用（資結書第~6.2~節）。

\textbf{㈢} 沿用題目的約定（Front 指向第一個元素的前一格、Rear 指向最後一個元素，開始時都是 $0$）。 Front $=$ Rear
表示空；為了和空分開，滿的條件是「Rear 再前進一格就碰到 Front」，所以留一格不用，最多存 N-1 個（要用滿 N 格可以另加
計數器或旗標）。
\begin{minted}[fontsize=\small]{cpp}
const int N = 10;
int Rear = 0, Front = 0;

void Insert(char item, char Queue[])
{
    if ((Rear + 1) % N == Front)          // 再放就和 Front 重疊：滿
        cout << "Queue Is Full";
    else {
        Rear = (Rear + 1) % N;            // 到 N-1 之後繞回 0
        Queue[Rear] = item;
    }
}

void Delete(char &item, char Queue[])
{
    if (Front == Rear)                    // 空
        cout << "Queue Is Empty";
    else {
        Front = (Front + 1) % N;
        item = Queue[Front];
    }
}
\end{minted}
兩個操作都是 $O(1)$。 例如 $N = 10$ 時插入 $9$ 個、刪除 $9$ 個之後 Front $=$ Rear $= 9$，再插入時 Rear 變成 $0$，
Queue[0] 也被用上，不會再有假溢位。 元素個數隨時是 $(\text{Rear} - \text{Front} + N) \bmod N$。
\end{pastsoln}
